%!TEX root = thesis.tex

\chapter{Evaluation}
\label{ch:evaluation}

This chapter presents the experimental setup, evaluation metrics, and results of the \ac{VL} segmentation framework.

\section{Experimental Setup}
\label{sec:eval:setup}

\subsection{Hardware and Software}
\label{sec:eval:hardware}

All experiments are conducted with the following configuration:
\begin{itemize}
	\item \textbf{GPU}: NVIDIA GPU with CUDA support.
	\item \textbf{Framework}: PyTorch 2.0+ with mixed-precision training (FP16).
	\item \textbf{Python}: 3.11 with conda environment management.
	\item \textbf{Key dependencies}: torchvision, transformers (HuggingFace), scikit-learn, scipy.
\end{itemize}

\subsection{Training Protocol}
\label{sec:eval:protocol}

Each model is trained for up to 100 epochs with early stopping (patience~15). The training set consists of 70 SUN cases, with 15 cases held out for validation and 15 for testing. Input images are resized to $352 \times 352$ pixels. The batch size is~8 for all experiments.

\subsection{Experimental Configurations}
\label{sec:eval:configs}

We evaluate six configurations as summarized in Table~\ref{tab:experiments}:

\begin{enumerate}
	\item \textbf{Vision-Only Baseline}: ResNet-34 encoder + U-Net decoder, no text input.
	\item \textbf{Full VL Model}: Complete pipeline with all four text attribute categories.
	\item \textbf{Shape Only}: Text restricted to Paris Classification morphology descriptions.
	\item \textbf{Size Only}: Text restricted to polyp size descriptions.
	\item \textbf{Location Only}: Text restricted to anatomical location descriptions.
	\item \textbf{Pathology Only}: Text restricted to surface pattern and boundary descriptions.
\end{enumerate}


\section{Evaluation Metrics}
\label{sec:eval:metrics}

\subsection{Segmentation Quality}
\label{sec:eval:seg_metrics}

\begin{itemize}
	\item \textbf{\ac{DSC}}: Measures overlap between predicted and ground truth masks:
	\begin{equation}
		\text{DSC} = \frac{2|P \cap G|}{|P| + |G|}
	\end{equation}
	
	\item \textbf{\ac{IoU}} (Jaccard Index): A stricter overlap measure:
	\begin{equation}
		\text{IoU} = \frac{|P \cap G|}{|P \cup G|}
	\end{equation}
	
	\item \textbf{Precision}: Fraction of predicted polyp pixels that are correct: $\frac{TP}{TP + FP}$.
	
	\item \textbf{Recall}: Fraction of actual polyp pixels correctly identified: $\frac{TP}{TP + FN}$.
\end{itemize}

\subsection{Uncertainty Quality}
\label{sec:eval:unc_metrics}

\begin{itemize}
	\item \textbf{\ac{URS}}: Quantifies uncertainty reduction from language integration (Section~\ref{sec:approach:urs}).
	\item \textbf{\ac{SAS}}: Measures alignment between text similarity and uncertainty patterns (Section~\ref{sec:approach:sas}).
	\item \textbf{\ac{AC-ECE}}: Attribute-conditioned calibration error (Section~\ref{sec:approach:acece}).
	\item \textbf{Mean Predictive Entropy}: Average entropy across test pixels.
	\item \textbf{Mean Predictive Variance}: Average MC~Dropout variance.
\end{itemize}


\section{Results on SUN Dataset}
\label{sec:eval:sun_results}

\subsection{Segmentation Performance}
\label{sec:eval:seg_results}

Table~\ref{tab:sun_results} presents the segmentation results across all six configurations on the SUN test set.

\begin{table}[htbp]
	\caption{Segmentation results on the SUN test set. Best results are highlighted in \textbf{bold}. Results to be populated after experimental runs.}
	\label{tab:sun_results}
	\centering
	\begin{tabular}{lcccc}
		\toprule
		\textbf{Configuration} & \textbf{Dice} $\uparrow$ & \textbf{IoU} $\uparrow$ & \textbf{Precision} $\uparrow$ & \textbf{Recall} $\uparrow$ \\
		\midrule
		Vision-Only      & ---    & ---    & ---    & ---    \\
		Full VL          & ---    & ---    & ---    & ---    \\
		Shape Only       & ---    & ---    & ---    & ---    \\
		Size Only        & ---    & ---    & ---    & ---    \\
		Location Only    & ---    & ---    & ---    & ---    \\
		Pathology Only   & ---    & ---    & ---    & ---    \\
		\bottomrule
	\end{tabular}
\end{table}

\subsection{Uncertainty Analysis}
\label{sec:eval:unc_results}

Table~\ref{tab:uncertainty_results} summarizes the uncertainty metrics. The \ac{URS} and \ac{SAS} require paired evaluation between the vision-only baseline and each \ac{VL} configuration.

\begin{table}[htbp]
	\caption{Uncertainty metrics on the SUN test set. URS and SAS are computed relative to the vision-only baseline. Results to be populated after experimental runs.}
	\label{tab:uncertainty_results}
	\centering
	\begin{tabular}{lcccc}
		\toprule
		\textbf{Configuration} & \textbf{Mean Entropy} $\downarrow$ & \textbf{Mean Var.} $\downarrow$ & \textbf{URS} $\uparrow$ & \textbf{SAS} $\uparrow$ \\
		\midrule
		Vision-Only      & ---    & ---    & ---   & ---   \\
		Full VL          & ---    & ---    & ---   & ---   \\
		Shape Only       & ---    & ---    & ---   & ---   \\
		Size Only        & ---    & ---    & ---   & ---   \\
		Location Only    & ---    & ---    & ---   & ---   \\
		Pathology Only   & ---    & ---    & ---   & ---   \\
		\bottomrule
	\end{tabular}
\end{table}


\subsection{Spatial Uncertainty Analysis}
\label{sec:eval:spatial}

Table~\ref{tab:spatial_urs} presents region-level URS values, comparing uncertainty reduction across polyp interior, boundary, and background regions.

\begin{table}[htbp]
	\caption{Region-level Uncertainty Reduction Score (URS). Positive values indicate language-driven uncertainty reduction. Results to be populated.}
	\label{tab:spatial_urs}
	\centering
	\begin{tabular}{lccc}
		\toprule
		\textbf{Configuration} & \textbf{Polyp URS} & \textbf{Boundary URS} & \textbf{Background URS} \\
		\midrule
		Full VL        & ---   & ---   & ---   \\
		Shape Only     & ---   & ---   & ---   \\
		Size Only      & ---   & ---   & ---   \\
		Location Only  & ---   & ---   & ---   \\
		Pathology Only & ---   & ---   & ---   \\
		\bottomrule
	\end{tabular}
\end{table}


\subsection{Attribute-Conditioned Calibration}
\label{sec:eval:acece_results}

Table~\ref{tab:acece} presents the \ac{AC-ECE} results conditioned on Paris Classification morphology classes.

\begin{table}[htbp]
	\caption{Attribute-Conditioned ECE by Paris Classification type. Lower values indicate better calibration. Results to be populated.}
	\label{tab:acece}
	\centering
	\begin{tabular}{lccccc}
		\toprule
		\textbf{Configuration} & \textbf{Is} & \textbf{Ip} & \textbf{IIa} & \textbf{IIb} & \textbf{Overall} \\
		\midrule
		Vision-Only  & ---   & ---   & ---    & ---    & ---   \\
		Full VL      & ---   & ---   & ---    & ---    & ---   \\
		\bottomrule
	\end{tabular}
\end{table}


\section{Ablation Study}
\label{sec:eval:ablation}

The ablation study isolates the contribution of each clinical text attribute to segmentation and uncertainty quality.

\subsection{Per-Attribute Contribution}
\label{sec:eval:perattr}

Table~\ref{tab:ablation_summary} summarizes the improvement of each single-attribute configuration over the vision-only baseline.

\begin{table}[htbp]
	\caption{Ablation study: improvement over vision-only baseline. $\Delta$Dice denotes the absolute improvement in Dice score. Results to be populated.}
	\label{tab:ablation_summary}
	\centering
	\begin{tabular}{lcccc}
		\toprule
		\textbf{Attribute} & $\Delta$\textbf{Dice} & $\Delta$\textbf{IoU} & \textbf{URS} & \textbf{SAS} \\
		\midrule
		Shape     & ---   & ---   & ---   & ---   \\
		Size      & ---   & ---   & ---   & ---   \\
		Location  & ---   & ---   & ---   & ---   \\
		Pathology & ---   & ---   & ---   & ---   \\
		All       & ---   & ---   & ---   & ---   \\
		\bottomrule
	\end{tabular}
\end{table}


\section{Cross-Dataset Generalization}
\label{sec:eval:crossdataset}

To assess generalization beyond the SUN training distribution, we evaluate the trained models on two external benchmarks without fine-tuning.

\begin{table}[htbp]
	\caption{Cross-dataset evaluation results. Models trained on SUN and evaluated on Kvasir-SEG and CVC-ClinicDB without fine-tuning. Results to be populated.}
	\label{tab:cross_dataset}
	\centering
	\begin{tabular}{lcccc}
		\toprule
		\multirow{2}{*}{\textbf{Configuration}} & \multicolumn{2}{c}{\textbf{Kvasir-SEG}} & \multicolumn{2}{c}{\textbf{CVC-ClinicDB}} \\
		\cmidrule(lr){2-3} \cmidrule(lr){4-5}
		& \textbf{Dice} & \textbf{IoU} & \textbf{Dice} & \textbf{IoU} \\
		\midrule
		Vision-Only & ---  & ---  & ---  & ---  \\
		Full VL     & ---  & ---  & ---  & ---  \\
		\bottomrule
	\end{tabular}
\end{table}

\section{Qualitative Results}
\label{sec:eval:qualitative}

\begin{figure}[htbp]
	\centering
	\fbox{\parbox{0.9\textwidth}{\centering\vspace{3cm}\textbf{[Qualitative Results Placeholder]}\\Columns: Input Image | Ground Truth | Vision-Only Prediction | Full VL Prediction | Uncertainty Map (Vision) | Uncertainty Map (VL)\vspace{3cm}}}
	\caption[Qualitative segmentation and uncertainty comparison.]{Qualitative comparison of segmentation predictions and uncertainty maps between the vision-only baseline and the full VL model. Each row shows a different test case. The uncertainty maps visualize predictive entropy, where brighter regions indicate higher uncertainty. To be populated after experimental runs.}
	\label{fig:qualitative}
\end{figure}

\section{Discussion}
\label{sec:eval:discussion}

The experimental results will be analyzed along the following dimensions once all runs are completed:

\begin{enumerate}
	\item \textbf{Segmentation improvement}: Whether the full VL model achieves statistically significant improvement over the vision-only baseline in Dice and IoU.
	
	\item \textbf{Attribute importance ranking}: Which individual attributes contribute most to segmentation quality and uncertainty reduction.
	
	\item \textbf{Uncertainty calibration}: Whether language guidance leads to better-calibrated predictions, particularly for difficult morphology types (e.g., flat lesions, IIb).
	
	\item \textbf{Spatial uncertainty patterns}: Whether URS is highest at polyp boundaries, indicating that language cues help most in ambiguous regions.
	
	\item \textbf{Generalization}: Whether VL models maintain their advantage on external benchmark datasets where no text metadata is available at test time (relying on default caption generation).
\end{enumerate}
